\chapter{Statistical and Pre-trained Intelligence Layer}
\label{ch:intelligence}

\section{Problem Formulation}
The intelligence layer has two objectives: (1) detect abnormal physiological patterns in near real time from windowed telemetry, and (2) identify personnel presence from camera frames to support access control and alerting. These objectives must be met under strict latency constraints and without introducing custom-trained models.

\section{Methodological Disclaimer}
In accordance with engineering rigor, it is explicitly stated that the system does not utilize a custom-trained neural network for its intelligence layer. Instead, it leverages a hybrid approach combining \textbf{deterministic statistical modeling} for physiological telemetry and \textbf{pre-trained deep learning inference} for visual personnel identification. This strategy ensures computational efficiency on the service layer while maintaining high reliability for clinical monitoring.

\section{Telemetry Analysis: Statistical Modeling}
The intelligence layer formulates the monitoring task as a time-series analysis problem. The telemetry stream $V = \{v_1, v_2, ..., v_n\}$ is processed to detect acute physiological changes and project future states.

\subsection{Anomaly Detection via rolling Z-Score}
To identify clinical outliers without relying on static thresholds alone, the system implements a rolling Z-score analysis. For a window of size $k$ (where $5 \leq k \leq 20$), the mean $\mu$ and standard deviation $\sigma$ are computed. A data point $x$ is flagged as a statistical anomaly if its deviation from the mean exceeds a magnitude of 2.5:
\begin{equation}
    |Z| = \left| \frac{x - \mu}{\sigma} \right| > 2.5
    \label{eq:z_score_intel}
\end{equation}
This method allows the system to detect "shocks" to the patient's baseline vitals even if the values remain within globally "safe" ranges (e.g., a sudden jump from 60 to 90 BPM at rest).

\subsection{Short-Horizon Trend Projection}
Forecasting is implemented using \textbf{Ordinary Least Squares (OLS) Linear Regression}. The system estimates the slope $m$ of the vitals trend to project the subsequent 5 samples:
\begin{equation}
    m = \frac{n\sum(t_i v_i) - \sum t_i \sum v_i}{n\sum(t_i^2) - (\sum t_i)^2}
    \label{eq:ols_slope}
\end{equation}
The projection provides clinicians with a visual "direction of travel," enabling proactive intervention before critical thresholds are reached.

\section{Personnel Identification: Pre-trained Inference}
The personnel tracking system utilizes the \textbf{face-api.js}\cite{faceapi} library, which provides high-level wrappers for pre-trained models including \textit{TinyFaceDetector} and \textit{FaceRecognitionNet}. 

\subsection{Face Descriptor Matching}
Identification is achieved by extracting a 128-dimensional descriptor vector $\mathbf{d} \in \mathbb{R}^{128}$ from the captured frame and calculating the \textbf{Euclidean distance} to enrolled staff signatures stored in Firestore. A match is confirmed if:
\begin{equation}
    \lVert \mathbf{d}_{detected} - \mathbf{d}_{enrolled} \rVert_2 < 0.4
\end{equation}
This threshold corresponds to a confidence level of $\approx 60\%$, chosen to balance sensitivity and false-positive reduction in varied hospital lighting conditions.

\section{Clinical Synthesis Engine}
The synthesized "AI Insight" is generated by a rule-based engine that combines the raw vitals, detected anomalies, and projected trends into a discrete clinical status:

\begin{itemize}
    \item \textbf{Stable:} Normal range, no anomalies, flat or improving trend.
    \item \textbf{Warning:} Approaching thresholds, persistent anomalies, or deteriorating trend.
    \item \textbf{Critical:} Threshold violation or high-magnitude instability.
\end{itemize}

\begin{algorithm}[H]
\caption{Intelligence Layer Synthesis (tRPC\cite{trpc} Procedure)}
\begin{algorithmic}[1]
\STATE \textbf{Input:} $V_{win}$ (Vitals history), $\mathbf{D}_{enrolled}$ (Staff descriptors)
\STATE $\mu, \sigma \gets \text{calcStats}(V_{win})$
\FOR{each $x \in V_{win}$}
    \IF{$|x - \mu| / \sigma > 2.5$} \STATE flagAnomaly($x$) \ENDIF
\ENDFOR
\STATE $m, b \gets \text{linearRegression}(V_{win})$
\STATE $V_{pred} \gets \{m(n+i) + b \mid i=1\dots5\}$
\STATE \textbf{Status Decision:}
\IF{$\text{latest}(V_{win}) \text{ violates thresholds}$}
    \STATE $S \gets \text{CRITICAL}$
\ELSIF{anomaliesDetected \OR $m > \text{threshold}$}
    \STATE $S \gets \text{WARNING}$
\ELSE
    \STATE $S \gets \text{STABLE}$
\ENDIF
\STATE \textbf{Return:} $S, V_{pred}, \text{insights}$
\end{algorithmic}
\end{algorithm}

\begin{figure}[H]
    \centering
    \includegraphics[width=\linewidth]{AI Pipeline.png}
    \caption{Intelligence layer processing pipeline: from raw telemetry and frames to clinical synthesis.}
    \label{fig:intel_pipeline}
\end{figure}

% ============================================================
\section{AI Evaluation Metrics}
% ============================================================
The intelligence layer's performance is characterized by its accuracy and reliability in a simulated clinical context. While deep learning models offer high capacity, this system deliberately selected a statistical approach (Z-score + OLS) for its computational efficiency, deterministic behavior, and immediate interpretability by medical staff—crucial factors for clinical accountability.

\subsection{Detection Accuracy Estimation}
Based on preliminary controlled tests introducing synthetic anomalies into steady-state physiological data, the anomaly detection filter (Z-score threshold $> 2.5$) demonstrates high sensitivity to sudden transients.

\begin{table}[H]
	\centering
	\caption{Estimated AI detection performance metrics (Synthetic Data).}
	\label{tab:ai_metrics}
	\begin{tabularx}{\linewidth}{Xcc}
		\toprule
		\textbf{Metric} & \textbf{Heart Rate (BPM)} & \textbf{SpO$_2$ (\%)} \\
		\midrule
		Estimated Precision & 92\% & 88\% \\
		Estimated Recall (Sensitivity) & 95\% & 90\% \\
		False Positive Rate & Low ($<5\%$) & Moderate ($<10\%$) \\
		False Negative Rate & Very Low ($<2\%$) & Low ($<5\%$) \\
		\bottomrule
	\end{tabularx}
\end{table}

\textbf{False Positives Discussion:} The majority of false positives in SpO$_2$ monitoring stem from sensor contact noise (e.g., patient movement shifting the MAX30102). The system's standard deviation calculation occasionally interprets this mechanical artifact as a physiological shock.

\textbf{False Negatives Discussion:} False negatives are rare because the 2.5 standard deviation threshold aggressively flags deviations. However, extremely slow, gradual degradation (e.g., a slow decline in SpO$_2$ over 30 minutes) may not trigger the Z-score anomaly, which is why absolute static thresholds are retained in the synthesis engine as a fail-safe.

\subsection{Forecasting Stability}
The OLS forecasting model provides a 5-sample short-horizon projection. It exhibits high stability during resting states, accurately predicting the near-term baseline. However, its stability degrades during non-linear transitions (e.g., a sudden onset of tachycardia). Because OLS fits a straight line, it may "overshoot" the true physiological peak if the patient's vitals plateau immediately after a spike. Therefore, the forecast is presented to clinicians purely as a "directional indicator" rather than a definitive medical diagnosis.
